% !TeX program = xelatex
% ============================================================
% applied/test_program_and_methodology/ieee_829.tex
% IEEE 829-2008 / ISO/IEC/IEEE 29119-3 — Software Test Documentation
%
% English international-standard ANALOG of the ЕСПД «Программа и
% методика испытаний» (ГОСТ 19.301-79). Selectable per §2.3.4.2 of the
% HSE SE practice program for students who write in English.
% Self-contained, compile-ready, English only (no language branch).
% Fill the yellow \hseFill{...} placeholders.
% ============================================================
\documentclass[12pt,a4paper]{extarticle}
\newcommand{\hseDefaultMono}{Consolas}
\input{../common/fonts}
\usepackage[english]{babel}
\usepackage[a4paper,left=25mm,right=15mm,top=20mm,bottom=20mm]{geometry}
\usepackage{setspace}\onehalfspacing
\usepackage{indentfirst}\setlength{\parindent}{1.25cm}\setlength{\parskip}{6pt}
\usepackage{microtype}\usepackage{csquotes}\usepackage{amsmath}\usepackage{graphicx}
\usepackage{xcolor}\usepackage{float}\usepackage{array}\usepackage{booktabs}
\usepackage{longtable}\usepackage{tabularx}\usepackage{adjustbox}\usepackage{xurl}\urlstyle{same}
\usepackage{hyphenat}\usepackage{caption}\usepackage{enumitem}\setlist{nosep}
\usepackage{titlesec}\usepackage{hyperref}
\hypersetup{unicode=true,colorlinks=true,linkcolor=black,urlcolor=blue}
\input{../common/meta}    % \hseFill, \projectname, \hseProjectNameEn, \hseAuthorName, \hseGroup, \hseYear ...
\input{../common/commands}
\begin{document}\sloppy

% ── Column type: top-aligned paragraph column ───────────────
\newcolumntype{L}[1]{>{\raggedright\arraybackslash}p{#1}}

% ── Title block ─────────────────────────────────────────────
\begin{center}
{\large\bfseries Software Test Documentation}\\[2pt]
{IEEE 829-2008 / ISO/IEC/IEEE 29119-3}\\[10pt]
{\bfseries Master Test Plan and Test Cases}\\[16pt]
{\large System under test:\\[2pt]\bfseries \hseProjectNameEn}\\[4pt]
{\itshape (\projectname)}\\[18pt]
Author: \hseAuthorName\\[2pt]
Group: \hseGroup\\[2pt]
Supervisor: \hseFill{supervisor full name}\\[2pt]
Year: \hseYear
\end{center}

\vspace{6pt}
\hrule
\vspace{10pt}

% ============================================================
%  PART I. MASTER TEST PLAN
% ============================================================
\section*{Part I. Master Test Plan}
\addcontentsline{toc}{section}{Part I. Master Test Plan}

% ── 1. Introduction ─────────────────────────────────────────
\section{Introduction}
This document defines the test plan for the software system \enquote{\hseProjectNameEn} (\projectname). It follows the structure of IEEE~829-2008 and ISO/IEC/IEEE~29119-3, and describes the scope, approach, resources, and schedule of the intended testing activities.

\textbf{Test scope.} \hseFill{state which parts of the system are covered by this plan}

\textbf{Objectives.} The testing aims to verify that the system satisfies its functional and non-functional requirements defined in the Software Requirements Specification (SRS) and to expose defects before release. Specific objectives:
\begin{itemize}
\item confirm that every requirement marked as testable is verified by at least one test case;
\item detect deviations between the observed and the expected behaviour;
\item \hseFill{add project-specific testing objective}
\end{itemize}

% ── 2. Test items ───────────────────────────────────────────
\section{Test items}
The following items and their versions are subject to testing.

\begin{longtable}{L{2.5cm} L{6.5cm} L{2.2cm} L{3.0cm}}
\toprule
\textbf{Item ID} & \textbf{Test item} & \textbf{Version} & \textbf{Source / reference} \\
\midrule
\endhead
ITM-01 & \hseFill{component or module name} & \hseFill{ver} & \hseFill{repo / build tag} \\
ITM-02 & \hseFill{component or module name} & \hseFill{ver} & \hseFill{repo / build tag} \\
\bottomrule
\end{longtable}

% ── 3. Features to be tested ────────────────────────────────
\section{Features to be tested}
The features below are within the scope of this plan. Each is traced to a requirement identifier (SRS-F-xx).
\begin{itemize}
\item SRS-F-01 --- \hseFill{feature to be tested}
\item SRS-F-02 --- \hseFill{feature to be tested}
\item SRS-F-03 --- \hseFill{feature to be tested}
\end{itemize}

% ── 4. Features not to be tested ────────────────────────────
\section{Features not to be tested}
The following features are excluded from this test plan, together with the reason for exclusion.
\begin{itemize}
\item \hseFill{excluded feature} --- \hseFill{reason (out of scope / third-party / not implemented)}
\item \hseFill{excluded feature} --- \hseFill{reason}
\end{itemize}

% ── 5. Approach ─────────────────────────────────────────────
\section{Approach}
Testing is organised into the following levels.
\begin{itemize}
\item \textbf{Unit testing} --- individual functions and classes are verified in isolation. Tool: \hseFill{unit framework, e.g. pytest / JUnit}.
\item \textbf{Integration testing} --- interaction between modules and external services is verified. Tool: \hseFill{integration approach / tool}.
\item \textbf{System testing} --- the assembled system is verified against the requirements end-to-end. Tool: \hseFill{system / e2e tool}.
\end{itemize}

\textbf{Methods.} \hseFill{black-box / white-box / equivalence partitioning / boundary-value analysis} are applied, complemented by manual exploratory testing where automation is not cost-effective.

\textbf{Automation and reporting tools.} \hseFill{CI system, coverage tool, tracker}.

% ── 6. Item pass/fail criteria ──────────────────────────────
\section{Item pass/fail criteria}
A test item is considered \emph{passed} when:
\begin{itemize}
\item all test cases with severity \enquote{critical} and \enquote{major} have the status \enquote{Passed};
\item the automated test suite is green and code coverage is at least \hseFill{target \%};
\item no open defect of severity \enquote{critical} remains.
\end{itemize}
Otherwise the item is considered \emph{failed} and returned for correction.

% ── 7. Suspension criteria and resumption requirements ──────
\section{Suspension criteria and resumption requirements}
\textbf{Suspension.} Testing of an item is suspended if \hseFill{blocking condition, e.g. the build does not start or a critical defect blocks the main scenario}.

\textbf{Resumption.} Testing resumes once \hseFill{condition for resumption, e.g. the blocking defect is fixed and a new build is delivered}. The affected test cases are re-executed from the beginning.

% ── 8. Test deliverables ────────────────────────────────────
\section{Test deliverables}
The following artefacts are produced during testing:
\begin{itemize}
\item this Master Test Plan;
\item the test-case specification (Part II);
\item test logs and defect reports;
\item the requirements-to-tests traceability matrix;
\item \hseFill{other deliverable, e.g. coverage report}.
\end{itemize}

% ── 9. Test environment / environmental needs ───────────────
\section{Test environment}
The environment required to execute the tests is described below.

\begin{longtable}{L{4.5cm} L{10.0cm}}
\toprule
\textbf{Element} & \textbf{Specification} \\
\midrule
\endhead
Hardware & \hseFill{CPU / RAM / disk} \\
Operating system & \hseFill{OS and version} \\
Runtime / dependencies & \hseFill{language runtime, DBMS, services} \\
Test data & \hseFill{fixtures / seed data / anonymised dataset} \\
Network & \hseFill{connectivity, endpoints, mocks} \\
\bottomrule
\end{longtable}

% ── 10. Responsibilities and staffing ───────────────────────
\section{Responsibilities and staffing}
\begin{longtable}{L{4.5cm} L{5.0cm} L{5.0cm}}
\toprule
\textbf{Role} & \textbf{Person} & \textbf{Responsibility} \\
\midrule
\endhead
Test engineer & \hseAuthorName & Designs, executes and reports test cases \\
Developer & \hseFill{name} & Fixes reported defects \\
Reviewer & \hseFill{supervisor name} & Reviews test results \\
\bottomrule
\end{longtable}

% ── 11. Schedule ────────────────────────────────────────────
\section{Schedule}
\begin{longtable}{L{6.5cm} L{3.5cm} L{4.0cm}}
\toprule
\textbf{Activity} & \textbf{Start} & \textbf{End} \\
\midrule
\endhead
Test planning & \hseFill{date} & \hseFill{date} \\
Test design & \hseFill{date} & \hseFill{date} \\
Test execution & \hseFill{date} & \hseFill{date} \\
Reporting & \hseFill{date} & \hseFill{date} \\
\bottomrule
\end{longtable}

% ── 12. Risks and contingencies ─────────────────────────────
\section{Risks and contingencies}
\begin{longtable}{L{5.5cm} L{2.5cm} L{6.0cm}}
\toprule
\textbf{Risk} & \textbf{Impact} & \textbf{Contingency} \\
\midrule
\endhead
\hseFill{risk, e.g. unstable environment} & \hseFill{high/med/low} & \hseFill{mitigation action} \\
\hseFill{risk, e.g. schedule slippage} & \hseFill{high/med/low} & \hseFill{mitigation action} \\
\bottomrule
\end{longtable}

% ============================================================
%  PART II. TEST DESIGN AND TEST CASES
% ============================================================
\section*{Part II. Test Design and Test Cases}
\addcontentsline{toc}{section}{Part II. Test Design and Test Cases}

% ── Requirements traceability ───────────────────────────────
\section{Requirements under verification}
The functional requirements from the SRS that are verified by this test set are listed below; each is referenced by test cases in the next section.

\begin{longtable}{L{2.8cm} L{9.0cm} L{2.4cm}}
\toprule
\textbf{Req. ID} & \textbf{Requirement} & \textbf{Priority} \\
\midrule
\endhead
SRS-F-01 & \hseFill{functional requirement} & \hseFill{must/should} \\
SRS-F-02 & \hseFill{functional requirement} & \hseFill{must/should} \\
SRS-F-03 & \hseFill{functional requirement} & \hseFill{must/should} \\
\bottomrule
\end{longtable}

% ── Test cases ──────────────────────────────────────────────
\section{Test cases}
Each test case is identified by a \texttt{TC-xx} code and mapped to the requirement it verifies. Fill the preconditions, steps and expected result; record the status after execution (Passed / Failed / Blocked).

\begin{longtable}{L{1.7cm} L{2.0cm} L{3.0cm} L{3.6cm} L{3.4cm} L{1.6cm}}
\toprule
\textbf{TC ID} & \textbf{Req.} & \textbf{Preconditions} & \textbf{Steps} & \textbf{Expected result} & \textbf{Status} \\
\midrule
\endhead
\multicolumn{6}{r}{\emph{continued on next page}}\\
\endfoot
\bottomrule
\endlastfoot
TC-01 & SRS-F-01 & \hseFill{initial state} & \hseFill{ordered steps} & \hseFill{expected outcome} & \hseFill{P/F} \\
\addlinespace
TC-02 & SRS-F-02 & \hseFill{initial state} & \hseFill{ordered steps} & \hseFill{expected outcome} & \hseFill{P/F} \\
\addlinespace
TC-03 & SRS-F-03 & \hseFill{initial state} & \hseFill{ordered steps} & \hseFill{expected outcome} & \hseFill{P/F} \\
\end{longtable}

\end{document}
